\section*{Results}

\subsection*{Most released architectures cannot forbid impossible tensor predictions}

We first audited which released atomistic architectures could, even in principle, guarantee a symmetry-mandated zero for a parity-odd tensor. Fifteen released models were classified by reading their source at a pinned version or commit, recording the deciding line for each; where source inspection could not settle the question, the model was probed numerically rather than guessed (Table~\ref{tab:audit}). Six architectures attach parity labels to their internal features and can therefore make an odd-parity output vanish structurally on a centrosymmetric input. Nine cannot: three are rotation-equivariant only, three carry a single unlabelled Cartesian vector channel, and three are invariant models that cannot express an equivariant tensor output at all. The rotation-only group is not a fringe --- it contains EquiformerV2 and the eSCN family, which are among the most widely deployed backbones on large-scale materials benchmarks, as well as Equiformer v1 as released. The full audit, with the deciding source line for every row, is given in Supplementary Note~1 and Supplementary Table~1.

Two rows illustrate why measurement, not reading, has to settle such classifications. NequIP exposes a boolean named \texttt{parity}, yet its own documentation shows that disabling it does not remove parity labels from the features; the flag is not an SO(3) toggle, and our SO(3) arms are therefore built by relabelling irreducible representations directly. Conversely, GotenNet propagates its features through custom attention blocks that are not parity-typed, and we expected them to discard parity; a reflection test on the built model showed the opposite --- its vector output obeys the polar-vector transformation law under improper operations to machine precision ($1.8\times10^{-15}$, versus $1.9$ for the pseudovector law) --- so the row moved from undetermined to parity-aware on the basis of the measurement.

\begin{table}[t]
\centering
\small
\caption{Prevalence of feature parity across fifteen released architectures. Each model was classified from its released source at a pinned version or commit; two rows (ICTP, GotenNet) were additionally settled by numerical reflection tests at random initialisation. The consequence column states what the classification implies for a parity-odd tensor predicted on a centrosymmetric crystal. Deciding evidence for every row is quoted in Supplementary Table~1.}
\label{tab:audit}
\begin{tabular}{@{}l >{\raggedright\arraybackslash}p{3.7cm} >{\raggedright\arraybackslash}p{5.1cm}@{}}
\toprule
Class & Models & Consequence for a parity-odd tensor \\
\midrule
Parity-aware (6) & NequIP, Allegro, MACE, MatTen, ICTP, GotenNet & odd-parity output is structurally zero on a centrosymmetric input \\
SO(3)-only (3) & Equiformer v1, EquiformerV2, eSCN & nothing forces an odd-parity output to vanish \\
Vector-only (3) & PaiNN, TorchMD-Net (ET), ViSNet & no typed irreducible representations; a rank-3 head requires extra machinery \\
Invariant (3) & SchNet, DimeNet++, FAENet & cannot express an equivariant tensor output \\
\bottomrule
\end{tabular}
\end{table}

\subsection*{A label-free test with matched-pair architectures}

The experimental design is summarised in Fig.~\ref{fig:concept}. Three parity-aware architectures --- NequIP, Allegro, and MACE --- were each built in two arms through a single irreducible-representation helper: an O(3) arm carrying the natural spherical-harmonic parity $(-1)^{l}$ on every feature, and an SO(3) arm in which every feature is declared parity-even. The geometric content of the two arms is identical; only the parity bookkeeping differs, and with it the reflection equivariance that the feature algebra enforces. Architecture, hyperparameters, data, splits, and random seeds are shared exactly. EquiformerV2 joins as a fourth core, unmodified, because it has no parity labels to toggle; that absence is the reason for its inclusion. The full grid is 7 arms $\times$ 4 targets $\times$ 3 seeds $=$ 84 runs.

Before any training, every arm passed a verification gate probing an internal feature layer with a random rotation and a random improper operation: O(3) arms must satisfy both transformation laws, SO(3) arms must satisfy rotation and fail reflection. All arms behaved as required, and the SO(3) parameter count exceeds its O(3) counterpart in every configuration (ratios 1.00--1.40), so reduced capacity cannot explain any SO(3) deficit reported below (gate results in Supplementary Note~2 and Supplementary Table~2; per-configuration parameter counts in Supplementary Table~8).

Four targets separate two properties that a single experiment would confound: being a parity-odd target, and being evaluated where symmetry forbids a nonzero value. The QM9 internal energy $U_0$ (parity-even scalar) and the elasticity tensor (parity-even rank 4) carry no odd-parity constraint at all. The QM9 molecular dipole is parity-odd, but no molecule in QM9 has a point group that forces the dipole to vanish, so the constraint is never active. Only the piezoelectric tensor, evaluated on centrosymmetric crystals, is both parity-odd and symmetry-forbidden. The evaluation population comprises 2,000 centrosymmetric insulators from the Materials Project\cite{jain2013materials}, each re-verified with spglib\cite{togo2024spglib} before use, in two coordinate variants: idealized (snapped onto the exact space group) and raw (as relaxed by density functional theory).

\begin{figure}[t]
\centering
\includegraphics[width=\textwidth]{fig1_concept}
\caption{A label-free test of physical validity. A centrosymmetric crystal is passed to two networks that are identical except for the parity labels on their internal features. Neumann's principle forces the true piezoelectric tensor of such a crystal to vanish identically. The O(3)-equivariant network reproduces the zero structurally, at any parameter values; the SO(3)-equivariant network, exactly rotation-equivariant but parity-blind, is free to predict a tensor of the magnitude of a genuine piezoelectric response. Predicted magnitudes shown are the median over the evaluation population described in the text. \textcolor{red}{ATTENTION: schematic drawn programmatically as a placeholder; consider replacing with a designed graphic before submission.}}
\label{fig:concept}
\end{figure}

\subsection*{Parity labels decide whether impossible predictions occur}

Figure~\ref{fig:headline}a shows the distribution of predicted piezoelectric magnitudes $\lVert e \rVert_F$ on the idealized evaluation population, where the true value is exactly zero for every structure. The two arms are separated by six orders of magnitude. Across all three O(3) arms and all seeds, not one structure-run of the 18,000 exceeds the operating threshold of 0.01\,C\,m$^{-2}$: the O(3) medians sit at $3\times10^{-7}$ to $3\times10^{-6}$\,C\,m$^{-2}$, the float32 noise floor. The SO(3) arms exceed the same threshold for 89.5\% (NequIP), 91.0\% (Allegro), 90.8\% (MACE), and 95.7\% (EquiformerV2) of the population --- we call these false flags --- with median magnitudes of 0.44 to 0.96\,C\,m$^{-2}$. For calibration, 98.5\% of the real piezoelectric tensors in the training data exceed the threshold and their median magnitude is 0.51\,C\,m$^{-2}$: the typical SO(3) violation is the size of a genuine piezoelectric response, not a numerical residue.

The separation does not depend on the operating point. Over the two decades of threshold from $10^{-4}$ to $10^{-2}$\,C\,m$^{-2}$ the SO(3) false-flag fraction changes by at most 4 percentage points per core while the O(3) fraction never exceeds 0.0007 (Fig.~\ref{fig:headline}b); the full curves over 25 thresholds and both coordinate variants are given in Supplementary Note~3, Supplementary Fig.~1, and Supplementary Tables~3--4. Nor does it depend on the coordinate variant: the SO(3) false-flag fractions of the three matched cores are identical to four decimal places between idealized and raw coordinates, and EquiformerV2's changes by less than 0.001. The O(3) arms show a small nonzero raw-variant response traced to a single structure whose as-relaxed coordinates are genuinely non-centrosymmetric at tolerances below the selection tolerance --- the guarantee detected a data artifact rather than violating symmetry --- as quantified in Supplementary Note~4 and Supplementary Table~5.

Paired by structure, the separation is statistically overwhelming: one-sided Wilcoxon signed-rank tests on the 2,000 paired violations give $p<10^{-300}$ for every core and seed, with the O(3) violation smaller for 96--98\% of individual crystals (Supplementary Note~5 and Supplementary Table~6). The violation pattern is also not an idiosyncrasy of any one model. Per-structure violations correlate between independently seeded runs of the same architecture (Spearman $\rho = 0.77$--$0.83$), and independent architectures flag nearly the same crystals: the flagged sets of any two SO(3) cores overlap with Jaccard index 0.93--0.99. The failure is a property of the symmetry class of the model, not of a particular implementation.

\begin{figure}[t]
\centering
\includegraphics[width=\textwidth]{fig2_headline}
\caption{Parity labels separate possible from impossible predictions by six orders of magnitude. (a) Distribution of predicted piezoelectric magnitude $\lVert e \rVert_F$ on 2,000 spglib-verified centrosymmetric crystals (idealized coordinates), for the O(3) and SO(3) arms of each core; the true value is exactly zero for every structure. Boxes span the 25th to 75th percentile of the per-structure distribution (mean of per-seed percentiles over 3 seeds), whiskers the 5th to 95th, horizontal line the median, isolated dash the maximum. Percentages above each SO(3) box give the fraction of structures exceeding the threshold. The dotted line marks the operating threshold and the dashed line the median magnitude of real piezoelectric tensors in the training data. (b) False-flag fraction as a function of threshold (idealized coordinates), mean over 3 seeds; shaded bands show one standard deviation. Blue curves, O(3) arms; red curves, SO(3) arms; markers distinguish cores. EquiformerV2 has no O(3) arm.}
\label{fig:headline}
\end{figure}

\subsection*{The constraint costs nothing where it is inactive}

If parity labelling handicapped models on ordinary tasks, removing it could be a defensible trade. It does not. On the three targets where no symmetry-forbidden value is ever evaluated --- $U_0$, dipole, and elasticity --- the two arms are statistically indistinguishable: the between-arm difference in test mean absolute error (MAE) is at most 0.94 pooled seed standard deviations, with a sign that is not consistent across cores (Fig.~\ref{fig:accuracy}a; complete accuracy tables in Supplementary Note~6 and Supplementary Table~7).

The dipole is the decisive control. It is parity-odd, like the piezoelectric tensor, yet the arms are within seed noise of each other, because QM9 contains no molecule whose symmetry forces the dipole to vanish: the constraint an O(3) model enforces is never active on this dataset. The dipole therefore separates the two properties the piezoelectric experiment would otherwise confound --- an odd-parity target, and an evaluation where symmetry forbids a nonzero value. Only the second produces a difference between the arms.

Where the constraint is active, it helps rather than costs. On the piezoelectric tensor the O(3) arm has lower in-distribution test error than the SO(3) arm for all three matched cores, by 3.4 to 4.6 pooled seed standard deviations with a consistent sign (Fig.~\ref{fig:accuracy}b). EquiformerV2's in-distribution error is comparable to the matched cores, which matters for interpretation: its near-universal false flagging on centrosymmetric crystals coexists with ordinary regression competence on non-centrosymmetric ones.

\begin{figure}[t]
\centering
\includegraphics[width=\textwidth]{fig3_accuracy}
\caption{Removing parity labels costs nothing where the parity constraint is inactive, and costs accuracy where it is active. (a) Difference in test mean absolute error (MAE) between the SO(3) and O(3) arms, in units of the pooled per-seed standard deviation, for the three matched cores on all four targets. The shaded band marks one standard deviation, within which differences are indistinguishable from seed noise. (b) Piezoelectric test MAE for both arms of each core (mean over 3 seeds; error bars, one standard deviation). EquiformerV2 has no O(3) arm. MAE, mean absolute error.}
\label{fig:accuracy}
\end{figure}

\subsection*{Training cannot buy the zero}

The natural objection is that the SO(3) failure is a data problem: no centrosymmetric crystal appears in the piezoelectric training set, so the model was never told the answer. We tested this directly by retraining every SO(3) arm on the training set augmented with 1,000 fresh spglib-verified centrosymmetric insulators labelled with the exact zero tensor, drawn from six centrosymmetric space groups, with hyperparameters, validation, and test partitions unchanged.

The decisive measurement is the in-training control: predictions on the very crystals the model was trained, with explicit zero labels, to call zero (Fig.~\ref{fig:training}a). After training, the models still false-flag 89.5--92.2\% of those crystals. The failure is not underfitting --- the zero-labelled rows are fit better than the real-tensor rows in training MAE --- and it is not a data-coverage problem, since held-out false-flag fractions barely move whether or not the held-out crystal's space group appeared in training (changes of 0.00--0.03; full tables in Supplementary Note~7 and Supplementary Table~10). Gradient descent drives the predictions towards zero, halving the violation medians, and cannot arrive, because the required value is exact. Augmentation meanwhile improves the ordinary regression error on non-centrosymmetric test crystals for every core: it produces better regressors that still predict impossible values.

Weighting the zero-labelled rows more heavily in the loss does not change the conclusion (Fig.~\ref{fig:training}b). Up-weighting them one-hundred-fold reduces the false-flag fraction on the trained-on crystals only from 0.90 to 0.68, leaves 87\% of held-out crystals from the same space groups flagged, and never approaches the exact zero that the O(3) arms deliver at any weight, for free; the violation median shrinks roughly ninefold but remains more than four orders of magnitude above the O(3) floor (Supplementary Table~11).

A second proposed repair operates at test time: antisymmetrise the output over inversion, $T_{\mathrm{sym}}(x) = [T(x) - T(I\!\cdot\!x)]/2$. This removes the false flags of the exactly-equivariant SO(3) cores, but vacuously: on a centrosymmetric crystal the inverted structure is the same periodic solid up to a permutation, so any permutation-invariant model returns $T(I\!\cdot\!x)=T(x)$ and the symmetrised output is identically zero for any input whatsoever --- no parity information is involved. The fix also presupposes knowing the target's parity in advance, which is precisely the knowledge O(3) features carry, and it fails on EquiformerV2, which violates the permutation identity the construction relies on (Supplementary Note~8 and Supplementary Table~12).

\begin{figure}[t]
\centering
\includegraphics[width=\textwidth]{fig4_training}
\caption{Zero-labelled training data and loss reweighting reduce, but cannot eliminate, impossible predictions. (a) False-flag fraction on the 1,000 centrosymmetric crystals that each SO(3) arm was trained on with explicit zero labels, after augmented retraining (mean over 3 seeds). The blue line marks the O(3) arms, which were not retrained and predict zero structurally. (b) False-flag fraction of the NequIP SO(3) arm as the loss weight on zero-labelled rows increases from 1 to 100, evaluated on the trained-on crystals and on held-out centrosymmetric crystals from space groups that were seen or unseen during training (mean over 3 seeds; error bars, one standard deviation).}
\label{fig:training}
\end{figure}

\subsection*{The guarantee tracks the physics of symmetry breaking}

Three further measurements locate the failure precisely in the missing parity, rather than in any generic deficiency of SO(3) models.

First, the guarantee switches off exactly when the symmetry does. Displacing rutile TiO$_2$ continuously along a polar mode ($\delta$ interpolating from the centrosymmetric parent to beyond the physical distortion, every frame symmetry-verified) shows the O(3) arms rising smoothly from machine zero ($1.9\times10^{-7}$ to $9.1\times10^{-7}$\,C\,m$^{-2}$ at $\delta=0$) as the polar distortion grows, while the SO(3) arms sit on a spurious floor of 0.09--0.28\,C\,m$^{-2}$ that is already present at $\delta=0$ and dominates the physical signal throughout the small-distortion regime (Fig.~\ref{fig:mechanism}a). Rutile is the load-bearing choice: its rotation subgroup 422 permits a rank-3 invariant, so only parity forbids a response at $\delta=0$.

Second, the guarantee is differentiable. For an O(3) model the identity $T(I\!\cdot\!x) = -T(x)$ forces the Jacobian $J = \partial T/\partial \mathbf{r}$ at any centrosymmetric geometry to anticommute with inversion, so its response is supported entirely on inversion-odd (polar) displacement modes. Measured over 20 crystals spanning 19 space groups and all seeds, with every Jacobian certified against finite differences, the even-subspace energy fraction of $J$ is $1.4\times10^{-7}$ to $9.7\times10^{-7}$ for the O(3) arms and 0.42--0.54 for the SO(3) arms --- five to six orders of magnitude apart with no overlap over any structure or seed (Fig.~\ref{fig:mechanism}b). The O(3) architecture does not merely output a zero; it encodes a physically structured map of which atomic displacements can activate a piezoelectric response.

Third, the crystals that SO(3) models do predict correctly are exactly the ones where rotations alone forbid the tensor. Splitting the evaluation population by point group (Fig.~\ref{fig:mechanism}c): on crystals of point group m$\bar 3$m, whose proper-rotation subgroup 432 admits no rank-3 invariant, the three matched SO(3) cores false-flag exactly none of the 166 crystals in any seed; on the neighbouring cubic class m$\bar 3$, whose rotation subgroup 23 does admit one, they flag 89--100\%. Both group-theoretic statements are verified numerically by Reynolds-operator rank tests in the released code. On cubic perovskites in the distortion experiment above, whose parent rotation subgroup is likewise 432, both arms start at machine zero and the sweep cannot distinguish them --- the same effect seen in a different guise. Of the whole population, 91.7\% lies outside what rotation can enforce --- closely matching the observed false-flag fractions. SO(3) models keep the zeros that rotations already guarantee; parity labels extend the guarantee to the strictly larger set that inversion demands (per-family detail in Supplementary Note~9 and Supplementary Table~14).

The same pattern holds for familiar named materials: on corundum Al$_2$O$_3$ (sapphire), a textbook centrosymmetric crystal whose non-cubic point group leaves parity as the only protection, the SO(3) arms predict magnitudes of 0.21 to 1.13\,C\,m$^{-2}$ against an exact zero, while rocksalt and diamond-structure compounds of class m$\bar 3$m are correctly near zero for all matched cores (Supplementary Table~13).

\begin{figure}[t]
\centering
\includegraphics[width=\textwidth]{fig5_mechanism}
\caption{The missing guarantee is specifically parity, and it is differentiable. (a) Predicted piezoelectric magnitude of rutile TiO$_2$ along a continuous polar distortion $\delta$ of the centrosymmetric parent (seed-averaged per core; open markers at the left edge show $\delta=0$). Blue, O(3) arms; red, SO(3) arms; markers distinguish cores; EquiformerV2 values are means over five evaluation draws. (b) Even-subspace energy fraction of the Jacobian $J=\partial T/\partial \mathbf{r}$ at centrosymmetric geometries of 20 crystals spanning 19 space groups, for both arms of the three matched cores (all seeds; each point one structure--seed pair). (c) False-flag fraction of the SO(3) arms split by point-group family of the evaluated crystal: m$\bar 3$m, where the proper-rotation subgroup alone forbids a rank-3 tensor; m$\bar 3$, where it does not; and all non-cubic classes. Bars annotated 0.000 are exactly zero in every seed. The blue line marks the O(3) arms, which are zero in every family.}
\label{fig:mechanism}
\end{figure}
